\begin{afterword}
\summaryhead

The post answers Perry Metzger's claim that as long as there are limited resources and competing actors that pass on their characteristics, there is selection pressure. The author asks how much. By Price's equation, a characteristic changes according to its covariance with relative fitness, and only to the extent that it is inherited. So a gene that saves its whole species does not spread.

Corporations compete, have offspring and die, but a spinoff resembles its parent only weakly, and there have been few corporate generations. Failed firms are weeded out, as short-lived stars vanish from the sky, but one-off weeding does not build complex features; those come from human design. Nanodevices whose instructions are encrypted, so that any copying error destroys them, would have almost no heritable variation, and would use up their resources within a few generations. Grey goo that ate the Earth would then barely change. Cumulative selection needs replication, variation in traits and in reproduction, their correlation, faithful heredity, frequent turnover, and all of this over many iterations.

\responsehead

The post's central distinction is sound and clearly made. Weeding out the unfit once, as with firms that go bankrupt or stars that burn out, is different from cumulative selection, which needs faithful inheritance over many generations to build complex features. The star example makes this plain, and the closing list adds useful quantitative conditions to the standard ones.

The thesis is worded more broadly than the argument. The post says that limited resources and competing actors that pass on characteristics are ``not sufficient to give rise to an evolution.'' On the standard account (Lewontin, 1970), variation, differential fitness and heritable fitness are enough for a population to ``undergo evolutionary change,'' so, given variation among the competitors, the quoted claim that there is ``selection pressure'' is true. What the post shows is that the pressure may be too weak or too short-lived to build complex adaptations. That is the question it asks, ``How much selection pressure?'', and the closing list answers it.

The case of corporations measures heredity from a parent firm to its spinoff. Evolutionary economists measure it differently. Nelson and Winter (1982) treat a firm's routines as its genes, inherited within the firm over time and gaining ground as the firms that have them grow. The post does not discuss this work, and its bit counts for corporate design and selection (``hundreds,'' ``thousands,'' ``a handful'') are not calculated.

One technical definition is wrong. Correlation is defined as covariance divided by variance; it is covariance divided by the product of the standard deviations. The paragraph's point still holds after the correction.

The other facts are correct. The Foresight guidelines say what the post reports, the resource limit holds (fewer than 270 doublings would use every atom in the observable universe), and the hedged account of George Price's death matches Wikipedia's.

In short: a sound distinction between one-off and cumulative selection, stated at first more broadly than the argument supports, and applied to corporations without discussing the work that treats firms' routines as inherited.
\end{afterword}
